% New English translation from checked cote 25, site 82 (PDF page 83).
% AI-translated by GPT-6 Astra (OpenAI, ChatGPT Pro mode).
\scanpage{25/82}
\mspage{\hand{77}}

as general base. Let us mention the following (without proof): let $f:X\hand{\to} S$
be a projective flat morphism, \var{\xout{\unsure{over S such that X}}}{}{typewriter strike-out}{ega25_25_082-v0001} with geometric fibers that are
\var{\xout{geome}}{}{typewriter strike-out}{ega25_25_082-v0002} irreducible and $(R_2)$, $S'$ a sub\var{\inserted[typed]{pre}}{pre}{insertion}{ega25_25_082-v0003}scheme of $X$ finite over $S$,
$s\hand{\in} S$, and suppose that for every $x\hand{\in} S'$ over $s$, $X_s$ is smooth over
$k(s)$ at $x$. Then there exist \unsure{an} open neighborhood $U$ of $s$, and an open
neighborhood $V$ of $S'|U$ in $X|U$, such that $V\hand{\longrightarrow} U$ is poly-elementary.
If $Y$ is a closed subprescheme of $X$\var{\inserted[typed]{not meeting $S'$ and}}{ not meeting $S'$ and}{insertion}{ega25_25_082-v0004} such that the set $Z$ of points
at which $Y$ is not smooth over $S$ \var{\xout{has fibers of}}{}{typewriter strike-out}{ega25_25_082-v0005} \var{\deleted{is such that}}{}{deletion}{ega25_25_082-v0006} \var{\inserted{satisfies}}{satisfies}{insertion}{ega25_25_082-v0007}
dim $Z_s$ \hand{$\leq$} dim $X_s$ - 2, then above we can take $V$ contained in
$X$ - $Y$.

c) (Possibly to be included in the statement of 13.2.) With
the notation of 13.2, let $\hand{\unsure{\phi|\emptyset}}$ be a finite subset of $X$ consisting of smooth
points of $X$, and suppose that $\hand{\unsure{\phi|\emptyset}}$ is contained in an \var{\xout{\illegible{1}}}{}{typewriter strike-out}{ega25_25_082-v0008} affine open subset
of $X$. Then $\hand{\unsure{\phi|\emptyset}}$ has a fundamental system of poly-elementary open
neighborhoods. \var{\inserted[typed]{Clearly we can assume $\hand{\unsure{\phi|\emptyset}}$ consists of \uline{closed} points.}}{Clearly we can assume $\hand{\unsure{\phi|\emptyset}}$ consists of \uline{closed} points.}{insertion}{ega25_25_082-v0009} The proof is essentially the same, except
that we formulate 13.3 in a slightly different form, namely: there
exists a linear subvariety $C$ of $P^r$ of codimension $n$, not \unsure{meeting}
$\hand{\unsure{\phi|\emptyset}}\hand{\cup} Y$, and such that for every $x_i\hand{\in}\hand{\unsure{\phi|\emptyset}}$\var{\inserted[typed]{$(\overset{\hand{-}}{k})$}}{$(\overset{\hand{-}}{k})$}{insertion}{ega25_25_082-v0010}, the linear
subvariety $L_i$ of codimension $n$-1 spanned by $C$ and $x_i$ has properties
a), b), c) of 13.3. To verify this point, we note that it suffices to
verify that the \uline{generic} $C$ has the stated properties; now \var{\inserted{for such a $C$}}{for such a $C$}{insertion}{ega25_25_082-v0011} \var{\deleted{then}}{}{deletion}{ega25_25_082-v0012}
\var{\xout{\unsure{the $L$}}}{}{typewriter strike-out}{ega25_25_082-v0013}\unsure{\hand{,}|$_i$} each $L_i$ is \unsure{generic} among the $L$ of codimension $n$ through
$x_i$, so that we can apply 13.3 in its initial form
(or at least in the form in which we proved it, which was: every
\uline{sufficiently general} $L_o$ through $x$ has properties a), b), c)).

e) One reason for the interest of 13.2 is the particularly simple
“topological” structure of elementary \var{\xout{\unsure{open subsets}}}{}{typewriter strike-out}{ega25_25_082-v0014} \var{\inserted[typed]{alg. schemes}}{alg. schemes}{insertion}{ega25_25_082-v0015}\var{\inserted{$U$}}{ $U$}{insertion}{ega25_25_082-v0016}. For example,
\tnote{25-082-1}{The repeated symbol alternatives phi/empty set are retained literally. The letter-or-stroke doubt within rencontrant is marked on meeting. The codimension n-1 of the first L with subscript i, and the later codimension n, are retained as written; no adjustment is made to the cited form of 13.3.}

\mspage{\hand{82}}
